\section{A compactly supported stretching step}
\label{step:section}

For a metric $h$ and a smooth function $t$, put
\[
 \mathcal D(h,t)=\Scal_h+2\lap_h t-|\dd t|_h^2.
\]
The next theorem enlarges the distance from a compact set while preserving
a quadratic reserve in this quantity.  Its error has two parts: a trace
decrement, which telescopes under successive metric enlargements, and an
arbitrarily small quadratic loss.

\begin{theorem}[Local stretching]\label{step:main}
Fix $n\ge2$ and $0<\delta<1/10$, and put $L_*=1.004$.  There are
constants $W,T\ge1$, depending only on $n,\delta$, with the following
property.  Let $h_0$ be a complete smooth metric on a manifold $M^n$
without boundary, let $h\ge h_0$ be smooth, and let $t\ge T$ be smooth.
Let $A$ be a nonempty compact set and $B$ a nonempty closed set such that
\[
 \dist_{h_0}(A,B)\ge W.
\]
For every $0<d<\delta$ there are smooth $h'\ge h$ and $t'\ge t$ such that
the changes have compact support in $\{\dist_{h_0}(A,\cdot)<W\}$,
\begin{equation}\label{step:conclusion}
 \mathcal D(h',t')\ge \mathcal D(h,t)+(t'^2-t^2)
 -\delta\bigl(\tr_h h_0-\tr_{h'}h_0\bigr)t'^2-dt'^2,
\end{equation}
and
\begin{equation}\label{step:distance}
 \dist_{h'}(A,B)\ge L_*W.
\end{equation}
In particular, $h'$ is complete.  Neither $W$ nor $T$ depends on $d$,
the geometry of $M$, or the sets $A,B$.
\end{theorem}

The proof uses a modified graph equation.  Its modification near zero
slope supplies an adjustable small error, whereas at moderate and large
slopes it is the ordinary graph equation.  A positive floor in the spatial
cutoff permits a closed-manifold existence argument.  We subsequently send
that floor to zero only on a fixed annulus, where the solution decays
exponentially.  All geometric estimates used for these two operations
take place on one compact auxiliary manifold.

\subsection{The slope switch and the order of constants}

We first specify the deformation associated with a provisional smooth
height $u$; the height itself will be chosen below. Write $s=|\dd u|_h$.
Fix $s_0=10^{-2}$ and a smooth nondecreasing function $\rho:[0,\infty)
\to[0,1]$ such that
\[
 \rho=0\text{ on }[0,s_0],\qquad
 \rho=1\text{ on }[5s_0,\infty),\qquad 0\le\rho'\le60.
\]
For $0<a\le1/2$ and $s\ge0$ define
\begin{equation}\label{step:switch}
 \mu=a^2+(1-a^2)\rho(s),\quad
 \theta=\frac{\mu s^2}{1+s^2},\quad
 l=(1-\theta)^{1/2},\quad \beta=\log l.
\end{equation}
For $s>0$ also put
\begin{equation}\label{step:chi}
 \chi=\left(1+\frac{a^2s_0^2}{\theta}\right)^{1/2},\qquad
 m=\frac{\dd\log\chi}{\dd\log\sqrt\theta}
   =-\frac{a^2s_0^2}{\theta+a^2s_0^2}.
\end{equation}
Here $\theta$ is strictly increasing.  The following bounds will be used
repeatedly:
\begin{equation}\label{step:switch-bounds}
\begin{gathered}
 0\le -s\beta_s\le2,\qquad -1\le m\le0,\qquad
 1\le\chi\le2\quad(s\ge s_0),\\
 m\le-\tfrac12,\quad \chi\ge s_0/s,\quad
 \theta\chi^2\le2a^2s_0^2\quad(0<s\le s_0).
\end{gathered}
\end{equation}
Indeed, on $s\le5s_0$,
\[
 -s\beta_s=\frac{s\theta_s}{2(1-\theta)}
 \le30s^3+\frac{s^2}{1+s^2}<2,
\]
and on $s\ge5s_0$ it equals $s^2/(1+s^2)$.  The remaining bounds follow
directly from $\theta\ge a^2s^2/(1+s^2)$ and the formula with $\mu=a^2$
on $s\le s_0$.

For $t\ge1$ we also have
\begin{equation}\label{step:beta-budget}
 \frac{(t-\beta)^2-t^2}{\theta}\le4(1+s)t\qquad(s>0).
\end{equation}
To check this, write $b_*=-\beta$.  If $\theta\le1/2$, then
$b_*\le\theta$, and the left side is at most $3t$.  Otherwise use
$b_*\le\tfrac12\log(1+s^2)\le\log(1+s)\le\sqrt s$ to bound it by
$4t\sqrt s+2s\le4(1+s)t$.

For the provisional height $u$, define the raw metric and drift by
\begin{equation}\label{step:raw-output}
 h_u=h+\frac{\theta}{(1-\theta)s^2}\,\dd u^2
     =h+\frac{\mu}{1+(1-\mu)s^2}\,\dd u^2,
 \qquad t_u=t-\beta.
\end{equation}
These expressions are smooth also at $s=0$, and $h_u\ge h$, $t_u\ge t$.
To see explicitly the smoothness needed at zero slope, note that near $s=0$
the coefficient of $\dd u^2$ is
$a^2/[1+(1-a^2)s^2]$, while the drift increment is
$-\tfrac12\log([1+(1-a^2)s^2]/[1+s^2])$.
Both are smooth functions of the first jet of $u$; therefore
$C^3$ convergence of $u$ to zero gives $C^2$ convergence of these
metric and drift increments to zero.
For $s\ge5s_0$, $h_u=h+\dd u^2$.  At points of nonzero slope in the original manifold, let
$e=\grad_hu/s$. With the given background metric $h_0$, the rank-one
inverse formula gives
\begin{equation}\label{step:trace}
 \tr_hh_0-\tr_{h_u}h_0=\theta b^2,
 \qquad b^2=h_0(e,e)\le1.
\end{equation}

We now fix all constants that are required to be independent of the local
geometry and of $d$.  Let
\begin{equation}\label{step:constants}
 A_0=1000,\qquad P_0=20,\qquad C_1=s_0^{-1},\qquad
 b_0=\frac{\delta s_0^2}{1000},\qquad
 0<\eps\le\min\left\{\frac{s_0}{100},\frac1{8n},\frac1{8C_1}\right\}.
\end{equation}
Apply \cref{pressure:law} with these constants, obtaining $F,S_0,T$.
In particular $S_0\le1$ and $T\ge1$.  Put $C_2=2C_1b_0$.
Only after a loss $d$ has been specified do we choose $a$ so that
\begin{equation}\label{step:a-choice}
 0<a\le\tfrac12,\qquad (2C_2+3)a^2\le d/2.
\end{equation}
The width $W$ will be a sufficiently large numerical constant; it may,
and will, be at least $10^5$.  It is chosen to satisfy the explicit
cutoff and distance inequalities below.  Thus the choices of $W,T$
precede $d,a$ and all local geometric choices.

\subsection{Cutoffs and a closed auxiliary manifold}

Write $r_A(x)=\dist_{h_0}(A,x)$.  There exist smooth functions
$\psi\ge0$ and $0\le\eta_1,\zeta\le1$ such that
\begin{equation}\label{step:cutoffs}
\begin{gathered}
 \inf_A\psi>W/5,\qquad |\dd\psi|_{h_0}\le1,
 \qquad\supp\psi\subset\{r_A<W/2\},\\
 \eta_1=1\text{ on }\{r_A\le .55W\},\qquad
 \eta_1=0\text{ on }\{r_A\ge .8W\},\qquad
 |\dd\eta_1|_{h_0}\le s_0/8,\\
 \zeta=1\text{ on }\{r_A\le .82W\},\qquad
 \zeta=0\text{ on }\{r_A\ge .98W\}.
\end{gathered}
\end{equation}
Choose also a fixed compact annulus $\mathcal A_\zeta$ in
$\{.8W<r_A<W\}$ containing $\overline{\{0<\zeta<1\}}$.
All cutoff error estimates below are imposed on this whole annulus,
including any intermediate plateau of $\zeta$.

Here no injectivity-radius or curvature bound is needed.  For clarity, a
smooth approximation to a Lipschitz function on the compact neighborhood
in question can be made in finitely many coordinate charts.  Shrink each
chart until its metric distortion is as small as desired, convolve there,
and combine with a smooth partition of unity.  The gradient of the
combination is the convex combination of the local gradients plus terms
bounded by the chosen uniform approximation error times the derivatives
of the partition.  Choose the local errors small enough to make the latter
sum arbitrarily small.  One may preserve prescribed constant regions by
using charts lying inside them and leaving a buffer between each such
region and the transition.  Apply this construction to
$(W/3-r_A)_+$, with value error less than $W/100$ and gradient bound
$1.05$, then divide by $1.05$ and use the zero-region buffer to obtain
$\psi$.  Smooth cutoff profiles of $r_A/W$ give the other two functions;
for $\eta_1$ the gradient can be bounded by $12/W<s_0/8$.  This proves
\cref{step:cutoffs} after increasing $W$ once.

For $0<r\le1/2$ define the positive-floor cutoff
\begin{equation}\label{step:floor}
 \eta=r+(1-r)\eta_1^2.
\end{equation}
Then $\eta=1$ near $\supp\psi$, $\eta=r$ on $\{r_A\ge.8W\}$,
and, uniformly in $r$,
\begin{equation}\label{step:floor-bounds}
 0<\eta\le1,\qquad |\dd\eta|_h\le s_0/4,
 \qquad \frac{|\dd\eta|_h^2}{\eta}\le1.
\end{equation}
The same statements hold with $h_0$ in place of $h$; this follows from
$h\ge h_0$ and $|\dd\eta|^2/\eta\le4|\dd\eta_1|^2$.

Completeness of $h_0$ and compactness of $A$ make the closed neighborhoods
used above compact.  Enclose a slightly larger neighborhood, say
$\{r_A\le1.1W\}$, in a smooth relatively compact domain.  Outside this
neighborhood, modify $h$ and $t$ near the boundary of the domain to be
products, and double the domain.  Convex interpolation for $t$ preserves
$t\ge T$, and interpolation of positive definite tensors gives the
metric extension.  We thereby obtain a closed auxiliary manifold $N$
on which $h,t,\psi,\eta_1,\zeta$ are smooth, with the last three functions
extended by zero where appropriate.  The extensions of $h,t$ are fixed
independently of $r$.  If the relevant component of $M$ is already compact,
it can simply be used as $N$.  The construction can be made on finitely
many components if $A$ meets more than one.  From now until the final
cutoff, we work on $N$.  Constants allowed to depend on this compact
geometry and on $a$ will be denoted by $C$; they will never enter $W,T$.

\subsection{The modified graph equation}

For a covector $p$ with $s=|p|_h$, set
\[
 \mathcal A(p)=\frac{\sqrt\mu}{\sqrt{1+s^2}}p^\sharp,
 \qquad J(s)=\sqrt{\theta+a^2s_0^2}.
\]
The symbols $\mu,\theta,l,\beta,\chi,m$ will always be evaluated at
$s=|\dd u|_h$ when a function $u$ is under discussion.  The maps
$\mathcal A,J$ are smooth at $p=0$: there $\mu=a^2$ is constant.
The tangential and radial eigenvalues of $\mathcal A_p$ are
$\sqrt\theta/s$ and $(\sqrt\theta)_s$, respectively; both are positive,
and both have limit $a$ at the origin.

For $\tau\in[0,1]$ consider
\begin{equation}\label{step:PDE}
 \Div_h\mathcal A(\dd u)-\dd t\bigl(\mathcal A(\dd u)\bigr)
 =J(|\dd u|_h)F\left(\frac{u-\tau\psi}{\eta},t\right).
\end{equation}
The pressure is evaluated only where its first argument lies in
$(-S_0,S_0)$.  We first prove existence with bounds independent of $r$.

\begin{lemma}[Height and gradient bounds]\label{step:existence}
For each fixed $a>0$ as above and every $0<r\le1/2$,
\cref{step:PDE} at $\tau=1$ has a smooth solution $u$ on $N$.
There are $S_1<S_0$ and $G<\infty$, independent of $r$ and of
$\tau\in[0,1]$, such that every solution along the continuation satisfies
\begin{equation}\label{step:C01}
 |u-\tau\psi|\le S_1\eta,\qquad |\dd u|_h\le G.
\end{equation}
The constants $S_1,G$ may depend on $a$ and on the compact data.
\end{lemma}

\begin{proof}
For $S<S_0$ the functions $u_\pm=\tau\psi\pm S\eta$ have uniformly
bounded first and second derivatives, independently of $r,\tau,S$.
Consequently the left side of \cref{step:PDE} at these functions is
bounded.  Since $J\ge as_0>0$ and $t(N)$ is compact, the uniform pole
limits of $F$ permit one choice $S_1<S_0$ for which $u_+$ has strictly
negative residual (left side minus right side) and $u_-$ has strictly
positive residual.  Comparison at a maximum or minimum, using strict
monotonicity in $u$, gives the height estimate.

One can make this argument before a solution has been constructed by
extending $F$ smoothly to all first arguments with $F_x>0$, retaining its
original values on a slightly larger interval than $[-S_1,S_1]$.  For
example, extend the positive even function $F_x$ by a smooth cutoff to a
positive constant and integrate from zero.  Solve this globally defined
equation.  The same barriers force every solution into $[-S_1,S_1]$, so
that the extension is never used in the resulting solution.

We give the gradient estimate, including the dependence of the pressure
on the graph angle.  If $\max_N|\dd u|_h\le1$ there is nothing to prove.
Otherwise, near a point of maximum slope, $\mu=1$ and the graph in
$(N,h)\times\R$ is an ordinary graph.  Write
\[
 w=(1+|\dd u|_h^2)^{-1/2},\qquad
 \nu=(N_u,-w),\qquad N_u=w\grad_hu,
\]
for its downward unit normal.  Denote its second fundamental form and
mean curvature by $\mathrm{II}$ and $H$, with $H=\Div\nu$.  Extend
$t$ constantly in the vertical direction and put $X=-\grad t$ and
$K=-\partial_z$.  The equation is
\[
 H+\langle X,\nu\rangle=p(x,z,w),\qquad
 p(x,z,\omega)=\sqrt{1-\omega^2+a^2s_0^2}\,
 F\left(\frac{z-\tau\psi(x)}{\eta(x)},t(x)\right).
\]
Let $K^\top=K-w\nu$.  Under a normal variation with speed $v$,
\begin{align*}
 \dot H&=-\lap v-(|\mathrm{II}|^2+\Ric(\nu,\nu))v,\\
 \frac{\dd}{\dd q}\langle X,\nu\rangle
 &=-\Hess t(\nu,\nu)v-\langle X^\top,\grad v\rangle,\\
 \dot w&=-\langle K^\top,\grad v\rangle.
\end{align*}
Thus the linearization of the equation, with all terms moved to the
left, is $-\mathcal Lv$, where
\begin{equation}\label{step:angle-operator}
 \mathcal L=\lap+
 \langle X^\top-p_\omega K^\top,\grad\,\cdot\,\rangle
 +|\mathrm{II}|^2+\Ric(\nu,\nu)+\Hess t(\nu,\nu)
 +\partial_\nu^{\mathrm{fixed}\ \omega}p.
\end{equation}
The last derivative differentiates the base and height variables while
holding the angle fixed.  Translation in direction $K$ leaves the angle,
$H$, and $\langle X,\nu\rangle$ unchanged and differentiates the
residual by $p_z$.  Its normal speed is $w$; its tangential component
differentiates an identically zero residual.  We obtain the exact identity
\begin{equation}\label{step:angle-Jacobi}
 \mathcal Lw=-p_z.
\end{equation}
In particular the term $-p_\omega K^\top$ in the drift has not been
discarded.

Put $\xi=(u-\tau\psi)/\eta$.  On the graph,
\[
 p_z=\frac{JF_x}{\eta},\qquad
 \partial_\nu^{\mathrm{fixed}\ \omega}p
 =J\left\{
 \frac{F_x}{\eta}\bigl[-w-\tau\dd\psi(N_u)-\xi\dd\eta(N_u)\bigr]
 +F_t\dd t(N_u)\right\}.
\]
The height bound puts $F,F_x,F_t$ in a fixed compact range.  Moreover
$|N_u|\le1$, $|\dd\psi|_h\le1$, and $|\dd\eta|_h\le s_0/4$.
The zero-order coefficient in \cref{step:angle-operator}, after omitting
the nonnegative term $|\mathrm{II}|^2$, is therefore at least $-C/\eta$,
with $C$ independent of $r,\tau$.  At the minimum of $w$ under
consideration, the slope is at least one, so $J\ge1/\sqrt2$ and
$p_z\ge A_0T/(\sqrt2\eta)$.  Since $\grad w=0$ and $\lap w\ge0$,
\cref{step:angle-Jacobi} implies
\[
 \frac{A_0T}{\sqrt2\eta}\le\frac C\eta w.
\]
This gives a positive lower bound for $w$, and hence the asserted $G$.

We finish the fixed-$r$ continuation argument in enough detail to specify
its regularity inputs. A classical $C^{2,\alpha}$ solution is smooth before
we apply estimates requiring smoothness. Indeed, writing its equation as
$a^{ij}u_{ij}=f$, the coefficient functions $a^{ij}$ and $f$ are
$C^{1,\alpha}$. On nested coordinate balls, a difference quotient $u_h$
satisfies
\[
 a^{ij}(x)\partial_{ij}u_h
 =f_h-(a^{ij})_h\partial_{ij}u(x+h).
\]
For each nonzero $h$, $u_h$ is already $C^{2,\alpha}$; its right side has
a $C^{0,\alpha}$ bound independent of $h$. The classical interior Schauder
estimate and passage to $h=0$ give $u\in C^{3,\alpha}$, and iteration gives
smoothness. These qualitative estimates may depend on the existing
$C^{2,\alpha}$ norm; the following a priori estimates do not.  At $\tau=0$, $u=0$ solves the equation.  For fixed
$r>0$, the height and gradient bounds put $(x,u,\dd u)$ in a compact set
on which the principal coefficients are uniformly elliptic and all
explicit coefficient derivatives are bounded.  In a coordinate chart
write the divergence-form equation as
\[
 \partial_i\mathfrak A^i(x,Du)=\mathfrak B(x,u,Du).
\]
Differentiating with $v=\partial_ku$ gives
\begin{equation}\label{step:differentiated}
 \partial_i(\mathfrak A^i_{p_j}\partial_jv)
 -\mathfrak B_{p_j}\partial_jv
 =\mathfrak B_u v+\mathfrak B_{x_k}
   -\partial_i\mathfrak A^i_{x_k}.
\end{equation}
The $x$ derivatives on the right are explicit derivatives of the
coefficient functions.  All nondifferentiated terms and the vector field
inside the last divergence are bounded using only the compact jet range;
no second derivative estimate has been assumed.  The local Hölder
estimate for uniformly elliptic divergence-form equations, with bounded
lower-order coefficients and bounded divergence data, now gives a
uniform $C^{1,\alpha}$ bound for $u$ on a finite cover.  The equation in
nondivergence form then has Hölder coefficients and right side, so local
Schauder estimates give $C^{2,\alpha}$ bounds, followed by bounds of
every order.  The H\"older and interior Schauder estimates used here
are respectively \citet[Lecture~18, Theorem~3, and Lecture~12,
Theorem~2]{simonPDE}.
All their uses are on compact coordinate balls with a positive
ellipticity constant; their constants may depend on fixed $r$.

Fix the floor $r>0$, and choose $0<\alpha<1$ smaller, if necessary,
than the exponent in the preceding H\"older estimate.
The linearization in $u$ has strictly negative zero-order coefficient
$-JF_x/\eta$; angle and gradient derivatives affect only its principal
and first-order terms.  We verify that this linearization is an
isomorphism from $C^{2,\alpha}(N)$ to $C^{0,\alpha}(N)$.

First consider $\lap_h-1$.  For $f\in L^2(N)$, minimize on $H^1(N)$
the strictly convex coercive functional
\[
 \frac12\int_N\bigl(|\dd v|_h^2+v^2\bigr)
                 \,\dd\operatorname{vol}_h
 +\int_N fv\,\dd\operatorname{vol}_h.
\]
Its Euler equation is $(\lap_h-1)v=f$ in the weak sense.
Testing the homogeneous difference with itself proves uniqueness
among $H^1$ weak solutions.
For $f\in C^{0,\alpha}(N)$ choose smooth $f_j$ converging uniformly
to $f$, with uniformly bounded $C^{0,\alpha}$ norms, using coordinate
mollification and a finite partition of unity.  Weak interior
regularity makes the corresponding weak solutions $v_j$ smooth
\citep[Lecture~6, Theorem~1 and Corollary~1]{simonPDE}.
The maximum principle gives $\|v_j\|_{C^0}\le\|f_j\|_{C^0}$.
A finite interior cover and the Schauder estimate
\citep[Lecture~12, Theorem~2]{simonPDE} then bound $v_j$ uniformly
in $C^{2,\alpha}(N)$.  Passing to a subsequence in
$C^{2,\alpha'}(N)$, where $0<\alpha'<\alpha$, gives a solution for $f$.
The common $\alpha$-H\"older bound for the second derivatives passes
to the limit, so this solution belongs to $C^{2,\alpha}(N)$.
Weak uniqueness identifies it with the energy solution, and the
maximum principle also gives uniqueness in the classical class.

Let $L_q$, $0\le q\le1$, interpolate the coefficients of
$\lap_h-1$ and the linearization written in nondivergence form.
Uniform ellipticity and a uniform
strictly negative zero-order bound persist along this path.  The
maximum principle and the same finite-cover Schauder estimate give
\[
 \|v\|_{C^{2,\alpha}(N)}
 \le C\|L_qv\|_{C^{0,\alpha}(N)},
 \qquad 0\le q\le1.
\]
Invertibility is open by the Neumann series.  To see closedness,
suppose $q_j\to q$ through invertible parameters and, for a fixed
$f\in C^{0,\alpha}(N)$, put $v_j=L_{q_j}^{-1}f$.
The estimate above bounds $v_j$ and gives
\[
 \|v_j-v_k\|_{C^{2,\alpha}}
 \le C\|(L_{q_k}-L_{q_j})v_k\|_{C^{0,\alpha}}
 \longrightarrow0.
\]
Thus $v_j$ converges in $C^{2,\alpha}$ to a solution of
$L_qv=f$; injectivity follows from the same estimate.
This proves the required linearized isomorphism.

The implicit function theorem now gives openness of the nonlinear
continuation in $\tau$.  The preceding difference-quotient bootstrap
applies to the $C^{2,\alpha}$ solutions obtained by implicit inversion.
The uniform a priori $C^{2,\alpha}$ bound makes its higher derivative
bounds uniform along continuation at this fixed positive floor.
In particular these bounds give compactness in
$C^{2,\alpha}(N)$, so a convergent sequence of soluble parameters
has a solution at its limit.  The height bound keeps the pressure
argument in its fixed compact interval.  The soluble set is therefore
closed as well as open and contains $0$, hence equals $[0,1]$.
The local bootstrap makes the resulting solution smooth.
\end{proof}

\subsection{Decay on the region of constant floor}

\begin{lemma}[Locality]\label{step:locality}
Let $u_r$ be the solutions from \cref{step:existence} at $\tau=1$.
On any compact subset of the fixed open region where $\psi=0$ and
$\eta_1=0$, the functions $u_r$ tend to zero in every $C^k$ norm as
$r\downarrow0$.
\end{lemma}

\begin{proof}
Here $\eta=r$ and $|u_r/r|\le S_1$.  Expanding the left side of
\cref{step:PDE} and using $F(0,t)=0$ writes its equation as
\begin{equation}\label{step:mass-equation}
 a^{ij}(x)\partial_{ij}u_r+b^i(x)\partial_i u_r
 -\frac{c_r(x)}r u_r=0,
 \qquad
 c_r=J\int_0^1F_x(qu_r/r,t)\,\dd q\ge as_0A_0T>0.
\end{equation}
The coefficients $a^{ij}$ are uniformly elliptic and bounded, and $b^i$
are bounded, independently of $r$.  Indeed they are smooth functions of
$(x,Du_r)$ on the fixed compact jet range from \cref{step:C01}; the
explicit $x$ part of the flux vanishes at zero covector and is therefore
a bounded first-order term.  Also $c_r$ is bounded above independently
of $r$.  No derivatives of these coefficients are needed for the next
comparison.

Take a coordinate box compactly contained in the constant-floor region.
Let $\ell_f$ be its inward affine distance from face $f$ in these
coordinates and set
\[
 V_r=M_0\sum_f\exp(-\gamma\ell_f/\sqrt r),
\]
where $M_0$ dominates $|u_r|$ on the boundary for every $r$.
Let $\Lambda$ bound $a^{ij}(\partial_i\ell_f)(\partial_j\ell_f)$,
let $c_*>0$ be the lower bound for $c_r$, and let $B_*$ bound the
first-order contributions.  The operator in \cref{step:mass-equation}
applied to each summand of $V_r$ is at most that summand times
\[
 \frac{\Lambda\gamma^2-c_*}{r}+\frac{B_*\gamma}{\sqrt r}.
\]
Choose $\gamma>0$ with $\Lambda\gamma^2<c_*/2$.  For small $r$ this
quantity is negative.  Comparison with both $u_r$ and $-u_r$ proves,
on each smaller box, $|u_r|\le C\exp(-c/\sqrt r)$.  A finite covering
gives this estimate on any fixed compact subset.

To upgrade the estimate to derivatives, rescale a coordinate ball of
radius comparable to $r$ by
\[
 x=x_0+ry,\qquad u_r(x)=r v(y).
\]
The equation becomes
\begin{equation}\label{step:rescaled}
 \partial_{y_i}\mathfrak A^i(x_0+ry,Dv)
 =r\mathfrak B_0(x_0+ry,v,Dv),
\end{equation}
where $\mathfrak B_0$ includes the drift term and $JF(v,t)$ and has no
negative power of $r$.  We have $|v|\le S_1$ and $|Dv|\le G$.
The principal ellipticity is uniform.  Explicit $y$ derivatives of
$\mathfrak A$ have a factor $r$, explicit $y$ derivatives of
$r\mathfrak B_0$ have a factor $r^2$, and its $v,p$ derivatives have a
factor $r$.  Thus all explicit coefficient derivatives are bounded
uniformly on this fixed jet range. Apply the H\"older estimate to the
differentiated \Cref{step:differentiated}, followed by the
Schauder bootstrap on fixed balls in $y$. For every $k$ this gives
\[
 |D_y^kv|\le C_k,\qquad |D_x^ku_r|\le C_kr^{1-k},
\]
on smaller balls, with constants uniform in $x_0$ in a fixed compact
set and in small $r$.  Finally, interior interpolation between the
exponentially small $C^0$ bound and the polynomial bound for one higher
derivative proves $C^k$ convergence to zero for each fixed $k$.
For example, on fixed nested compact subsets of the constant-floor region,
interpolation bounds the third derivatives by a constant times
$r^{-9/4}\exp(-c/(4\sqrt r))+\exp(-c/\sqrt r)$, which tends to zero.  This
last case ensures $C^2$ convergence of all metric and drift increments
in \cref{step:raw-output}, which are smooth functions of $Du_r$.
\end{proof}

\subsection{The scalar identity and its two slope regimes}

All comparisons involving $h_0$ in this subsection are made on the
original neighborhood $\{r_A<1.1W\}$, where the auxiliary data agree
with the given data.  No extension of $h_0$ to the double is needed.

Apply the construction in \cref{step:raw-output} to the solution $u_r$
from \cref{step:existence} at $\tau=1$, and write $u=u_r$. We now prove
its scalar estimate before applying the final spatial cutoff.

\begin{lemma}[Raw scalar estimate]\label{step:raw-estimate}
On this original neighborhood, the raw output satisfies
\begin{equation}\label{step:raw-inequality}
 \mathcal D(h_u,t_u)\ge\mathcal D(h,t)+(t_u^2-t^2)
 -\delta(\tr_hh_0-\tr_{h_u}h_0)t_u^2-\frac d2t^2.
\end{equation}
\end{lemma}

\begin{proof}
We derive the scalar comparison explicitly.  Near a point with $\dd u\ne0$
use $q=u$ as the transverse coordinate and write
\[
 h=v^2\dd q^2+\gamma_q,\qquad v=1/s,\qquad
 e=v^{-1}\partial_q,\qquad
 \mathrm J=\frac1{2v}\partial_q\gamma_q,\qquad H=\tr\mathrm J.
\]
Here $\mathrm J$ is the second fundamental form of the level sets,
and all tangential norms and derivatives below use $\gamma_q$.
Put $S=H-et$.  The Gauss equation and normal variation of $H$ give
\[
 \Scal_h=\Scal_{\gamma}-|\mathrm J|^2-H^2-2eH
              -2\frac{\lap_{\gamma}v}{v},
 \qquad
 \lap_ht=\lap_{\gamma}t+e(et)+H et
           +\langle\dd_\gamma\log v,\dd_\gamma t\rangle.
\]
Consequently
\begin{align}
 \mathcal D(h,t)
 &=-|\mathrm J|^2-S^2-2eS+\Scal_\gamma
    +2\lap_\gamma t-|\dd_\gamma t|^2\notag\\
 &\hspace{8mm}-2\frac{\lap_\gamma v}{v}
    +2\langle\dd_\gamma\log v,\dd_\gamma t\rangle.
 \label{step:foliation-scalar}
\end{align}
The deformation replaces $v$ by $v/l$ and $t$ by $t-\beta$.
The change in the tangential terms on the right is
\[
 2\langle\dd_\gamma\log v,\dd_\gamma\beta\rangle
       -|\dd_\gamma\beta|^2
 =c_\beta(2-c_\beta)|\dd_\gamma\log v|^2\ge0,
 \qquad c_\beta=-s\beta_s\in[0,2].
\]
This computation follows by expanding
$\lap_\gamma(v/l)/(v/l)$ and $|\dd_\gamma(t-\beta)|^2$;
the $\lap_\gamma\beta$ and the terms involving $\dd_\gamma t$ cancel.

The divergence equation can be written
\[
 S+e\log\sqrt\theta=\chi F(\xi,t),\qquad
 \xi=(u-\psi)/\eta.
\]
Define
\[
 M=\chi F,\qquad K=\frac{l^2}{\theta}e\beta.
\]
Because $e\log\sqrt\theta=-K$, we have $S=M+K$.
For the new metric, $e'=le$, $\mathrm J'=l\mathrm J$, and
\[
 S'=l(S+e\beta)=lM+K/l.
\]
Substitution in the normal terms of \cref{step:foliation-scalar} gives
\begin{equation}\label{step:E}
 \mathcal D(h_u,t_u)-\mathcal D(h,t)\ge\theta E,
 \qquad E=|\mathrm J|^2+M^2+K^2/l^2-2MK+2eM.
\end{equation}
For completeness, the derivative terms before simplification are
$-2le(lM+K/l)+2e(M+K)$.  Their $eK$ terms cancel; using
$e\beta=\theta K/l^2$ leaves $2\theta eM-2\theta MK+
2\theta K^2/l^2$, which combines with the squared terms to yield
\cref{step:E}.

Since $e\chi=-m\chi K$ and $et=H-M-K$,
\[
 eM=-mKM+\chi(F_t et+F_xe\xi).
\]
Completing the $H$ and $K$ squares gives $E\ge\mathcal B$, where
\begin{equation}\label{step:B}
 \mathcal B=\chi^2\left\{
 [1-l^2(1+m)^2]F^2
 -2[1+l^2(1+m)]FF_t-nF_t^2\right\}
 +2\chi F_xe\xi.
\end{equation}
An exact verification, useful for tracking every sign, is
\begin{align*}
 E-\mathcal B
 &=\left(|\mathrm J|^2-\frac{H^2}{n-1}\right)
 +\frac{(H+(n-1)\chi F_t)^2}{n-1}\\
 &\quad+\left(\frac K l-l\chi[(1+m)F+F_t]\right)^2
 +\theta\chi^2F_t^2\ge0.
\end{align*}
All pressure expressions here are evaluated at $(\xi,t)$.

Suppose first that $s\ge s_0$.  Denote the expression in braces in
\cref{step:B} by $\mathcal C$.  Its $F^2$ coefficient is nonnegative,
and $1\le1+l^2(1+m)\le2$.  From \cref{pressure:estimates} we get
\begin{equation}\label{step:brace-estimates}
 \mathcal C\ge-5\eps F_x,
 \qquad
 -FF_t\ge P_0F_x\ \Longrightarrow\ \mathcal C\ge P_0F_x.
\end{equation}
Indeed, if $FF_t\ge0$, then $F_t^2\le\eps FF_t$ and
$\mathcal C\ge-(4+n\eps)FF_t\ge-5\eps F_x$.
If $FF_t<0$, the same absolute-value estimate gives
$\mathcal C\ge(2-n\eps)|FF_t|\ge|FF_t|$.

Where $\eta\ne1$ we have $\psi=0$, and
\[
 e\xi=\frac{s-\xi e\eta}{\eta}\ge\frac{3s}{4\eta}\ge\frac{3s}4.
\]
Using $1\le\chi\le2$, \cref{step:constants,step:brace-estimates}
give $E\ge sF_x$ there.  Where $\eta=1$, we have
$e\xi=s-e\psi\ge s-b$, with $b^2=h_0(e,e)$ as in \cref{step:trace},
because $|\dd\psi|_{h_0}\le1$.
If $s>8b$, the same estimates give
\[
 E\ge(7s/4-20\eps)F_x\ge sF_x.
\]
It remains to consider $s_0\le s\le8b$.  The possible deficit from
$sF_x$ is at most $10F_x$: indeed
$2\chi(s-b)-s\ge-4$ and $20\eps<1$.
On the favorable pressure branch, however,
\[
 E-sF_x\ge[20\chi^2+2\chi(s-b)-s]F_x\ge0,
\]
since $s\le8b\le8$.  Thus a deficit is possible only when
$F_x\le b_0t^2$.  In that case $b\ge s_0/8$ and our choice of $b_0$
gives $10b_0\le\delta(s_0/8)^2$.  We conclude throughout $s\ge s_0$
that
\begin{equation}\label{step:large-slope}
 E\ge sF_x-\delta b^2t^2
 \ge4(1+s)t-\delta b^2t^2.
\end{equation}
The last inequality follows from $F_x\ge1000t$ and $s\ge.01$.
Combining \cref{step:beta-budget,step:E,step:trace} proves the raw
inequality in this region, even without its $d$ term.

Now suppose $0<s<s_0$.  Since $m\le-1/2$, the first coefficient in
$\mathcal C$ is at least $3/4$.  Using $|F_t|\le\eps|F|$ gives
\[
 \mathcal C\ge(3/4-4\eps-n\eps^2)F^2\ge\tfrac12F^2.
\]
Where $\eta=1$, $e\xi/\chi\ge-1$.  Elsewhere, a negative $e\xi$
can occur only if $s\le|\dd\eta|_h$, and then
\[
 \frac{|e\xi|}{\chi}
 \le\frac{|\dd\eta|_h s}{s_0\eta}
 \le\frac{|\dd\eta|_h^2}{s_0\eta}\le C_1.
\]
By the quadratic pressure upper bound,
\begin{equation}\label{step:small-slope}
 E\ge\chi^2\{(1/2-2C_1\eps)F^2-2C_1b_0t^2\}
 \ge-C_2\chi^2t^2.
\end{equation}
Since $\theta\chi^2\le2a^2s_0^2$, the loss in \cref{step:E} is
at most $2C_2a^2t^2$.  Also $-\beta\le\theta\le a^2s_0^2$ and
$t\ge1$, so $t_u^2-t^2\le3a^2t^2$.  The choice in
\cref{step:a-choice} pays both terms with $(d/2)t^2$.
The trace penalty is favorable and may be omitted in this regime.
This proves \cref{step:raw-inequality} wherever $s>0$.
At a zero of $\dd u$ it follows by continuity unless $\dd u$ vanishes
on a neighborhood; on such a neighborhood the raw output equals the
input and the inequality is immediate.
\end{proof}

\subsection{Cutting off the deformation and obtaining distance}

Define on the original manifold
\begin{equation}\label{step:final-output}
 h'=h+\zeta(h_u-h),\qquad t'=t+\zeta(t_u-t),
\end{equation}
and extend by $(h,t)$ outside the auxiliary neighborhood.  The result
is smooth, $h'\ge h$, $t'\ge t$, and its changes have compact support
in $\{r_A<W\}$.  We now choose $r$ small enough to verify both
required estimates.

First, the scalar inequality must use the trace of the actual cutoff
metric.  On the fixed compact annulus $\mathcal A_\zeta$, define the residual for any positive definite metric $H$
and function $z$ by
\[
 \mathfrak E(H,z)=\mathcal D(H,z)-\mathcal D(h,t)-(z^2-t^2)
 +\delta\bigl(\tr_hh_0-\tr_Hh_0\bigr)z^2+dz^2.
\]
This is a continuous function of the $C^2$ jets of $(H,z)$ on a compact
set of positive definite metrics.  By \cref{step:raw-inequality},
\[
 \mathfrak E(h_u,t_u)\ge dt_u^2-(d/2)t^2\ge(d/2)T^2.
\]
The annulus lies compactly in the fixed region $\psi=0$, $\eta_1=0$.
By \cref{step:locality}, $u_r\to0$ there in $C^3$; hence both the raw
increments and their products with $\zeta$ tend to zero in $C^2$.
Choose $r$ so small that
\[
 |\mathfrak E(h',t')-\mathfrak E(h_u,t_u)|<dT^2/4
\]
on this annulus.  Then $\mathfrak E(h',t')>0$ there.  Off the annulus
the output is either raw or unchanged, and the same desired inequality
holds directly.  This proves \cref{step:conclusion}.  In particular,
no rank-one trace identity for the uncut metric has been substituted
for the trace of the cutoff metric.

It remains to prove the distance gain, including for curves that revisit
the cutoff annulus or leave the auxiliary neighborhood.  Introduce the
globally defined effective height
\[
 v_{\mathrm{eff}}=\zeta u,
\]
extended by zero outside the neighborhood.  On $A$ the height barrier
and \cref{step:cutoffs} give $v_{\mathrm{eff}}>W/5-1$, whereas on $B$
it is zero.  Decrease $r$ further, if necessary, so that
\begin{equation}\label{step:height-cutoff}
 |\dd v_{\mathrm{eff}}|_h\le .01
\end{equation}
throughout $\mathcal A_\zeta$; locality justifies this choice.
Consider a piecewise smooth curve from $A$ to $B$ of $h'$ length at most
$L_*W$.  Its $h$ length is at most $L_*W$.  Partition its parameter interval according to $0<\zeta<1$,
$\zeta=1$, and $\zeta=0$. At points where $\zeta$ equals either endpoint,
$\dd\zeta=0$ because $\zeta$ is smooth and valued in $[0,1]$.
The total variation of
$v_{\mathrm{eff}}$ on $\{0<\zeta<1\}$ is at most $.01L_*W$,
counting all repeated visits.  On the region where $\zeta=1$ and
$s\le5s_0=.05$, its variation is at most $.05L_*W$.  Where $\zeta=0$
it has zero variation.  Therefore the variation on
\[
 \Omega_{\mathrm{graph}}=\{\zeta=1,\ s\ge .05\}
\]
is greater than
\[
 (1/5-.06L_*)W-1>.11W.
\]
The strict last inequality holds for our fixed sufficiently large $W$.
On this region $h'=h+\dd u^2$, while everywhere $h'\ge h$.
If $\gamma$ is parametrized on an interval $I$, the integral triangle
inequality in $\R^2$ yields
\begin{align*}
 \operatorname{Length}_{h'}(\gamma)
 &\ge\int_I\sqrt{|\dot\gamma|_h^2+
      \mathbf1_{\Omega_{\mathrm{graph}}}|(u\circ\gamma)'|^2}
       \,\dd q\\
 &\ge\left\{
   \operatorname{Length}_h(\gamma)^2+
   \left(\int_{\gamma^{-1}(\Omega_{\mathrm{graph}})}
     |(u\circ\gamma)'|\,\dd q\right)^2
   \right\}^{1/2}\\
 &>W\sqrt{1+.11^2}>L_*W.
\end{align*}
Here $\operatorname{Length}_h(\gamma)\ge W$ follows from $h\ge h_0$
and the original separation of $A,B$.  This contradiction proves
\cref{step:distance}; curves with infinite length are irrelevant.
Finally, a smooth metric dominating a complete metric is complete:
an $h'$-Cauchy sequence is $h_0$-Cauchy, and local equivalence of the
two smooth metrics supplies convergence in $h'$.  The proof of
\cref{step:main} is complete.
